% \documentclass[pairquote, minted]{hzguide}
\documentclass[pairquote]{hzguide}

\setmainfont{Noto Serif CJK KR}[BoldFont={* Bold}]
\setsansfont{Noto Sans CJK KR}[BoldFont={* Medium}]

\IfLuatex{
    \hangulpunctuations=0 
    \setmonohangulfont{Noto Sans Mono CJK KR}[Scale=0.95]    
}[
    \disablekoreanfonts    
    \setmonofont{D2Coding}
]

\LayoutSetup{
    paper=A4, 
    column=vartwo,
    ulmargin=25mm,
    ulratio=1.5,
    hook={\setheadfoot{10mm}{15mm}},
    % showtrims,
    % showlayout
}

\DecolorHyperlinks[MidnightBlue]
\HeadingSetup{
    cftsection,
    chapterstyle=sparse,
    cjkchapter,
    chaptercolor=MidnightBlue, 
    sectioncolor=MidnightBlue, 
    sectionstyle={firm, rule},
    paragraphstyle=indent,
    linespacing=1.33,
}
\sloppy
\PageStyleSetup{
    headrule=false,
    headouter={},
    evenfootleft={\thepage\quad |\qquad \leftmark},
    oddfootright={\rightmark\qquad |\quad \thepage},
    sectionmark=true
}
\ThumbIndexSetup{
    xoffset=0mm, toffset=50mm, boffset=50mm, 
    width=3em, height=2em, interval=5mm, 
    fgcolor=white, bgcolor=darkgray, 
    style=\sffamily\bfseries\LARGE\centering, 
    content=\thechapter
}
\makeatletter
\patchcommand{\@smemmain}{}{\ThumbIndexEnable} % mainmatter*
\makeatother
\patchcommand{\printindex}{\cleartoverso\ThumbIndexDisable}{}
\ObjectSetup{
    beforeskip=.1\onelineskip, 
    afterskip=.1\onelineskip,
    align=\centering
}   
\CodeSetup{
    beforeskip=.5\onelineskip,
    afterskip=-.1\onelineskip
}
\BoundedBoxSetup{
    beforeskip=.75\onelineskip,
    afterskip=.1\onelineskip,
    style=\setSpacing{1.1}
}

\RenewTerms{macros}{
    labeltype=macro,
    delimiter=\hfill,
    labelwidth=4.75em,
    leftmargin=5.35em,
    style=\raggedright,
}
\IndexingEnable*

\ExplSyntaxOn
\NewDocumentCommand \keyvalue { v }
{
    \group_begin:
        \ttfamily =~#1
    \group_end:
    \newline
}
\NewDocumentCommand \keyvalueTF {}
{
    \texttt{=~true/false}\newline    
}
\NewDocumentCommand \true {}
{
    \texttt{true}
}
\NewDocumentCommand \false {}
{
    \texttt{false}
}

\NewDocumentCommand \hzcls { m m }
{
    \inputminted[linenos=true, firstline=#1, lastline=#2]{latex}{../../../texmf/tex/latex/HzGuide/hzguide.cls}
}
\ExplSyntaxOff
\AnnotateSetup{index, star={index=false}}

\CoverSetup{
    AfterFront=\cleartorecto,
    BeforeBack=\cleartoverso,
    title={HzGuide},
    subtitle={테크니컬 라이터를 위한 레이텍 클래스},
    ProductImage=cover_image.jpg,
    DocumentType={사용 설명서},
    PubYear=2021,
    revision={2.0},
    note={\hfill 미니멀리즘의 목표는 사용자에게 요구되는 수고를 최소화하는 것이다.},
    manufacturer={이호재\quad\email{yihoze@gmail.com}},
    address={서울 문정동에서}
}

\begin{document}

\frontmatter*

\FrontCover
% \begin{IfVartwoEnlarge}
\TableOfContents*
\ListOfFigures*
% \end{IfVartwoEnlarge}

\chapter{머리말}

\begin{flushright}
\url{https://github.com/YiHoze/HzGuide}
\end{flushright}

이 문서는 루아텍으로 만들어졌다.
지텍으로 컴파일하면 소스 코드를 부여주는 부분에서 결과가 달라질 것이다.
나는 루아텍을 선호한다. 
지텍에 비해 조금 느리지만, 폰트 처리에서 더 합리적이기 때문이다.

memoir 클래스와 expl3 문법에 친숙한 사용자는, 그런 사용자가 HzGuide 클래스를 사용할 이유가 있을지 모르겠지만, 이 문서를 이해하는 데에 어려움이 없을 것이다.

\mainmatter*

\chapter{개요}

\term{HzGuide}는 \term{memoir}에 기반하여 \term{테크니컬 라이터}\annotate{technical writer}의 설명서 작성을 돕기 위한 목적으로 만들어진 레이텍 클래스이다.

\section{클래스 옵션}

메므와 클래스 옵션들과 함께 다음 옵션들을 사용할 수 있다.

\begin{code}
\documentclass[language=korean, styleset=../foo.tex, Noto]{hzguide}
\end{code}

\begin{macros}<hzguide>
\item[language] \keyvalue{korean, english, chinese, czech, danish, dutch, finnish, german, italian, japanese, norwegian, polish, portugeue, russian, slovakian, spanish, swedish, turkish, TC}
\macro{chinese}는 \term{간체}\annotate{简体}, \macro{TC}는 \term{번체}\annotate{繁體}이다.
영어는 \verb{language=english} 대신 \verb{english}로 지정할 수 있다.
언어에 따라 \macro{polyglossia} 또는 \macro{xeCJK} 패키지가 호출된다.

\item[styleset] \keyvalue{foo.tex}
지정된 설정 파일을 불러들인다.

\item[Noto] \keyvalueTF
Noto 폰트가 사용된다.

\item[minted] \keyvalueTF
\macro{minted} 패키지가 사용된다.. \macro{\coderead} 명령이 \macro{\verbatiminput} 대신 \macro{\inputminted}를 사용하게 된다.

\item[pairquote] \keyvalueTF
\macro{csquotes} 패키지가 사용된다.

\item[template] \keyvalueTF
\macro{hztemplate.tex}을 불러들인다. 설명서 초안을 만들기 위한 목적으로 설명서에 사용되는, 아래와 같은, 여러 요소들을 예시하는 명령들을 제공한다.

\begin{code}
\chapter{Introduction}\tplpara
\section{Features}\tpllist
\section{Package Items}\tplimagetable
\section{Specifications}\tplspectables
\end{code}
\end{macros}

\ChapterContentsEnable

\chapter{페이지 모양과 문서 구조}

\section{판형}

\macro{\LayoutSetup} 명령을 이용하여 판면과 여백을 설정할 수 있다.

\begin{code}
\LayoutSetup{
    paper=A4, 
    column=vartwo
}
\end{code}

\begin{macros}<\LayoutSetup>
\item[paper] \keyvalue{A4, A3, B5, A5, A5V, letter, 46D, slide, arbitrary}
일반적인 규격에서 벗어나는 판형을 만들려면 \macro{arbitrary}를 지정하고, 아래 옵션들을 적절하게 설정하라.

\item[stockwidth] \keyvalue{250mm}
인쇄 용지의 폭. 재단선이나 반달 색인을 표시해야 할 때 페이지보다 넓은 크기가 지정되어야 한다.

\item[stockheight] \keyvalue{353mm}
인용 용지의 높이

\item[paperwidth] \keyvalue{210mm}
페이지 폭

\item[paperheight] \keyvalue{297mm}
페이지 크기

\item[landscape] \keyvalueTF
페이지가 가로 방향으로 바뀐다.

\item[column] \keyvalue{one, vartwo} 
\term{변이단}\annotate{變二段}으로 조판하려면 \macro{vartwo}를 지정하라.

\item[ulmargin] \keyvalue{36mm}
상단 여백

\item[ulratio] \keyvalue{1.0}
하단 여백을 결정할 배율을 지정하라. "2"를 지정하면 하단 여백이 상단 여백의 두 배가 된다.

\item[lrmagin] \keyvalue{35mm}
좌우 여백

\item[vartwomargin] \keyvalue{20mm}
변이단으로 조판할 때 오른쪽 여백. 왼쪽 여백은 오른쪽 여백의 세 배가 된다.

\item[showtrims] \keyvalueTF
재단선이 표시된다.
그림 \ref{trims_layout}\을 보라.

\item[showlayout] \keyvalueTF
이 옵션은 \macro{\ShowPageLayout} 명령을 호출하여 여백 영역과 텍스트 영역의 테두리에 선을 그린다.
그림 \ref{trims_layout}\을 보라.

\item[hook] 
\keyvalue{\setheadfoot{0mm}{5mm}}
이 예와 같은 구획 명령을 이 설정에 끼워넣을 수 있다.
\end{macros}

\image[float, scale=0.5]{trims_layout}(이 문서의 레이아웃)


\section{면주}

\macro{\PageStyleSetup} 명령을 이용하여 면주 모양을 변경할 수 있다. \macro{\LayoutSetup} 명령이 이 명령을 이용한다.

\begin{code}
\PageStyleSetup{
    headrule=true,
    evenheadleft=\leftmark,
    oddheadright=\rightmark,
    evenfootleft=\thepage,
    oddfootright=\thepage,
    chaptermark=true,
    sectionmark=true,
}
\end{code}

\begin{macros}<\PageStyleSetup>
\item[font] \keyvalue{\rmfamily\small}
폰트

\item[headleft] \keyvalue{\leftmark, \rightmark, \thepage, 텍스트}
상단 왼쪽 

\item[headcenter] 상단 가운데 
\item[headright] 상단 오른쪽
\item[headinner] 상단 안쪽 (왼쪽 페이지의 오른쪽과 오른쪽 페이지의 왼쪽)
\item[headouter] 상단 바깥쪽 (왼쪽 페이지의 왼쪽과 오른쪽 페이지의 오른쪽)
\item[footleft] 하단 왼쪽
\item[footcenter] 하단 가운데 
\item[footright] 하단 오른쪽 
\item[footinner] 하단 안쪽
\item[footouter] 하단 바깥쪽
\item[evenheadleft] 왼쪽 페이지의 상단 왼쪽
\item[evenheadcenter] 왼쪽 페이지의 상단 가운데
\item[evenheadright] 왼쪽 페이지의 상단 오른쪽
\item[oddheadleft] 오른쪽 페이지의 상단 왼쪽
\item[oddheadcenter] 오른쪽 페이지의 상단 가운데
\item[oddheadright] 오른쪽 페이지의 상단 오른쪽
\item[evenfootleft] 왼쪽 페이지의 하단 왼쪽
\item[evenfootcenter] 왼쪽 페이지의 하단 가운데
\item[evenfootright] 왼쪽 페이지의 하단 오른쪽
\item[oddfootleft] 오른쪽 페이지의 하단 왼쪽
\item[oddfootcenter] 오른쪽 페이지의 하단 가운데
\item[oddfootright] 오른쪽 페이지의 하단 오른쪽
\item[headrule] \keyvalueTF 
상단 가로선
\item[footrule] \keyvalueTF 
하단 가로선
\item[chaptermark] \keyvalueTF 
\macro{\chaptermark} 명령이 활성화된다.
\item[sectionmark] \keyvalueTF 
\macro{\sectionmark} 명령이 활성화된다.
\item[chapterpage] \keyvalueTF 
상단 면주가 장 페이지에도 적용된다.
\end{macros}

\section{장, 절, 문단}

\macro{\HeadingSetup} 명령을 이용하여 장과 절 제목의 모양을 그리고 문단 모양을 변경할 수 있다. 

\begin{code}
\HeadingSetup{
    linespacing=1.25,
    chapterfont=\normalfont\bfseries,
    sectioncolor=MidnightBlue, 
    paragraphstyle=indent 
}
\end{code}

\begin{macros}<\HeadingSetup>
\item[linespacing] \keyvalue{1.25} 
줄 간격. 이것은 \macro{\baselinestretch}를 설정하는 것과 같다.

\item[article] \keyvalueTF
\macro{article} 스타일이 적용된다. 이를테면 차례 제목이 절 수준으로 바뀐다.

\item[chapterwider] \keyvalueTF
변이단 조판에서 장 제목이 왼쪽 여백에서 시작되게 한다.

\item[chapterstyle] \keyvalue{tight, sparse}
장 스타일. 이것은 \macro{\chapterstyle}를 설정하는 것과 같다. 
그러나 메므와 클래스에 정의된 여러 장 스타일들 가운데 일부에 대해서만 아래 옵션들이 유효할 것이다.
\macro{tight} 스타일은 디폴트에서 줄 간격을 줄인 것에 불과하다.
\macro{veelo} 같은 스타일을 이용하려면, \macro{\printerchaptername}을 비롯한 여러 매크로들을 재정의해야 한다.

\item[cjkchapter] \keyvalueTF
한국어에 대해서는 "제\,1\,장"이, 중국어와 일본어에 대해서는 "第\,1\,章"이 식자된다.

\item[chapteralign] \keyvalue{\raggedleft, \centering, \raggedright}
장 제목의 정렬

\item[chapterfont] \keyvalue{\normalfont\bfseries}
장 제목의 폰트

\item[chaptercolor] \keyvalue{black}
장 제목의 색

\item[chapternamesize] \keyvalue{\huge}
장 이름의 크기 

\item[chapternumbersize] \keyvalue{\HUGE}
장 번호의 크기

\item[chaptertitlesize] \keyvalue{\Huge}
장 제목의 크기 

\item[nochaptername] \keyvalueTF 
장 이름이 식자되지 않는다. "제"도 식자되지 않는다.

\item[chaptercontents] \keyvalueTF
장 제목 아래에 장 차례가 만들어진다.

\item[sectionstyle] \keyvalue{firm, sparse, tight, multiline, rule}
절 스타일. \macro{multiline}은 변이단 조판에서 절 제목의 둘째 줄이 들여쓰기 되는 것을 방지한다.
두 가지 옵션을 함께, 이를테면 \macro{multiline}과 \macro{rule}을 동시에 지정할 수 있다.

\item[sectionalign] \keyvalue{\raggedleft, \centering, \raggedright}
절 제목의 정렬

\item[sectionfont] \keyvalue{\normalfont\bfseries}
절 제목의 폰트

\item[sectioncolor] \keyvalue{black}
절 제목의 색

\item[sectionsize] \keyvalue{\Large}
절 제목의 크기

\item[subsectionsize] \keyvalue{\large}
하위 절 제목의 크기

\item[subsubsectionsize] \keyvalue{\normalsize}
최하위 절 제목의 크기

\item[paragraphstyle] \keyvalue{interval/indent}
절 모양. 디폴트는 \macro{interval}이다. 이것은 들여쓰기 없이 문단 사이를 띄운다.
설정에 따라 목록들의 항목 간격도 함께 조정된다.

\item[parskip] \keyvalue{0.25\onelineskip}
문단 간격. \macro{paragraphstyle} 옵션에 \macro{interval}을 지정했을 때 이 값이 유효하다.

\item[parindent] \keyvalue{1em}
들여쓰기 크기. \macro{paragraphstyle} 옵션에 \macro{indent}를 지정했을 때 이 값이 유효하다.
\end{macros}

\section{장 차례}

\macro{\HeadingSetup}에 \macro{chaptercontents} 옵션을 지정하면 \macro{\ChapterContentsEnable} 명령이 호출된다.

\begin{code}
\ChapterContentsEnable[장 번호](절 수준)
\ChapterContentsDisable
\end{code}

장 번호 없이 \macro{\ChapterContentsEnable}을 선언하면 모든 장에 장 차례가 만들어진다.
장 번호를 지정하면 그 장에만 차례가 만들어진다.
절 수준에 "3"을 지정하면 최하위 절\annotate*{\string\subsubsection}까지 장 차례에 포함된다.
디폴트는 "2"\annotate*{\string\subsection}이다.
아래 예와 같이 \macro{\ChapterContentsDisable} 명령을 함께 사용하면 장 차례가 만들어져야 할 장의 범위를 지정할 수 있다.

\begin{code}
\ChapterContentsEnable
\chapter{troubleshooting}
... 
\chapter{Frequently Asked Questions}
...

\ChapterContentsDisable
\chapter{Warranty}
\end{code}

\section{사용자 행동}

하나의 \term{사용자 업무}\annotate{user task}는 하나 이상의 사용자 행동이 연결되어 달성된다.
소프트웨어 제품에서, 특히 모바일 앱에서, 매우 다양한 기능들이 일련의 단순한 조작으로 작동한다.
아래 예와 같은 설명은 "목적"과 "방법" 또는 "심리적 행동"과 "물리적 행동"으로 이루어져 있다.

\begin{example}
\vspace*{-1.5\onelineskip}
\paragraph{Turn iPhone on.} Press and hold the Sleep/Wake button until the Apple logo appears.
\paragraph{Unlock iPhone.} Press either the Sleep/Wake or Home button, then drag the slider.
\end{example}

\noindent 이와 같은 단락을 만드는 데에 \macro{\paragraph} 명령이 충분하지만, 자유롭게 스타일을 변경할 수 있도록 \macro{\action} 명령이 고안되었다.

\begin{code}
\ActionSetup{옵션}
\action{목적 텍스트}
\end{code}

\begin{macros}
\item[beforeskip] \keyvalue{0.5\hzparksip}
단락 위 간격

\item[afterskip] \keyvalue{-\hzparksip}
단락 아래 간격. \macro{inline} 옵션이 \true이면 이 간격이 무시된다.

\item[font] \keyvalue{\bfseries}
폰트 

\item[delimiter] \keyvalue{:, \quad}
구획 문자.

\item[inline] \keyvalueTF
이 옵션이 \false이면 방법 설명이 다음 줄에서 시작한다.
\end{macros}

이 단락들은 외형적으로 목록과 흡사하다.
하지만 이것들은, 각 단락의 \term{화제}\annotate{topic}가 서로 다르기 때문에, 목록이 될 수 없다. 

\section{워터마크}

\macro{\watermark} 명령을 이용하여 \term{워터마크}를 삽입할 수 있다.

\begin{code}
\watermark{
    x=15, y=150,
    fontsize=45, color=yellow, rotate=-90,
    text={WATERMARK}
}
\end{code}
\codeinput

워터마크를 없애려면 \macro{\watermark}가 선언된 페이지로부터 적어도 1 페이지 뒤에서 \macro{\ClearWatermark} 명령이 주어져야 한다.

\begin{code}
\WatermarkSetup{옵션}
\watermark{옵션}
\ClearWatermark
\end{code}

\begin{macros}
\item[x] \keyvalue{10}
x 좌표. 단위는 밀리미터이다.

\item[y] \keyvalue{10}
y 좌표

\item[color] \keyvalue{blue}
색

\item[font] \keyvalue{\sffamily}
폰트

\item[fontsize] \keyvalue{45}
폰트 크기. 단위는 포인트이다.

\item[rotate] \keyvalue{90}
기울기. 1 사분면에서 반시계 방향으로 주어진 각도만큼 회전한다.

\item[text] \keyvalue{텍스트}
워터마크 텍스트

\item[image] \keyvalue{이미지 파일}
\macro{text} 옵션과 \macro{image} 옵션이 동시에 설정된다면, 앞서 주어진 옵션이 무시된다.

\item[scale] \keyvalue{1.0}
주어진 배율만큼 이미지가 확대되거나 축소된다.
\end{macros}

\ClearWatermark

\section{반달 색인}

\macro{\ThumbIndexEnable} 명령을 이용하여 장 번호나 제목을 \term{반달 색인}\annotate{thumb index}에 넣을 수 있다. 반달 색인은 오른 페이지에만 생긴다.

\begin{code}
\ThumbIndexSetup{
    xoffset=0mm, toffset=40mm, boffset=40mm, 
    width=3em, height=2em, interval=3em, 
    fgcolor=white, bgcolor=darkgray, 
    style=\sffamily\bfseries\LARGE\centering, 
    content=\thechapter
}
\end{code}

이 기능은 \macro{\chaptermark} 명령을 이용하기 때문에 \verb{\chapter} 명령이 주어지지 않으면 발효되지 않는다.

\begin{code}
\ThumbIndexSetup{옵션}
\ThumbIndexEnable
\ThumbIndexDisable
\end{code}

위치나 크기 옵션에 밀리미터 단위로 값을 지정하는 것이 좋다.
주어진 값들로부터 밀리미터 단위로 변환하여 좌표를 구하기 때문에 그리고 이 과정에서 소수점 이하가 버려지기 때문에 정확한 계산을 기대하기 어렵다.

\begin{macros}
\item[xoffset] \keyvalue{0mm}
반달 색인이 바깥쪽으로 돌출되는 크기. 그림 \ref{trims_layout}\을 보라.

\item[toffset] \keyvalue{30mm}
첫 장에서 반달 색인의 y 좌표. 엄밀히 말해 이것은 페이지 상단으로부터 띄어야 하는 최소 거리이다.

\item[boffset] \keyvalue{30mm}
페이지 하단으로부터 띄어야 하는 최소 거리

\item[interval] \keyvalue{5mm}
앞 장의 반달 색인과 뒷 장의 반달 색인 사이의 간격

\item[width] \keyvalue{3em}
박스 폭

\item[height] \keyvalue{2em}
박스 높이 

\item[fgcolor] \keyvalue{white}
글자 색

\item[bgcolor] \keyvalue{gray}
박스 색

\item[style] \keyvalue{\bfseries\centering}
폰트와 정렬

\item[vertical] \keyvalueTF
박스가 세로 방향으로 회전한다.

\item[content] \keyvalue{\thechapter}
장 제목을 표시하는 방법에 대해 아래 예를 보라.
\end{macros}

\begin{code}
\ThumbIndexSetup{
    xoffset=12.1em, toffset=30mm, boffset=60mm, 
    width=15em, height=2.5em, interval=5mm, 
    fgcolor=white, bgcolor=darkgray, 
    style=\sffamily\bfseries\Large, 
    content=\quad\leftmark,
    vertical=true
}
\end{code}

\chapter{그림과 표}

\noindent 다른 종류의 문서들에 비해 설명서가 갖는 두드러진 특징들 중 하나는 그림과 표를 많이 포함한다는 것이고 그래서 그것들의 배치를 텍에 맡겨 놓기 곤란하다는 것이다.

\section{문단 사이에 이미지}

\macro{\image} 명령이 그림을 문단 사이에 둔다.

\begin{coderesult}
\image[scale=0.5]{uncertain}
(초고에서 그림 자리를 보여주기 위한 이미지)<초고 이미지>
\end{coderesult}

\begin{code}
\ImageSetup{옵션}
\image*|[옵션]{이미지 파일}[다른 이미지 파일](캡션)<짧은 캡션>
\end{code}

일반 출판물과 달리 설명서에서 그림은 보충적인 자료가 아니라 핵심적인 정보의 일부이기 때문에 \term{캡션}이 오히려 군더더기이다. 불가피한 경우에만 그림에 캡션을 붙여야 한다.

\begin{macros}<\ImageSetup>
\item[beforeskip] \keyvalue{\hznulldim, 0.5\onelineskip}
그림 위 간격. 값이 \macro{\hznulldim}이면 \macro{\ObjectSetup}으로 설정된 값이 사용된다.

\item[afterskip] \keyvalue{\hznulldim, 0.5\onelineskip}
그림 아래 간격. \macro{\hznulldim}이면 \macro{\ObjectSetup}으로 설정된 값이 사용된다.

\item[gap] \keyvalue{1em}
다른 이미지 파일이 주어졌을 때 두 이미지 사이 간격

\item[align] \keyvalue{\raggedleft, \centering, \raggedright}
이미지 정렬

\item[float] \keyvalueTF
이 옵션은 이미지를 \macro{figure} 환경에 넣는다. 다시 말해 그림의 위치가 텍에 의해 결정된다.

\item[float-placement] \keyvalue{tbh}
figure 환경을 이용하는 경우에 위치 옵션

\item[frame] \keyvalueTF
이미지 테두리에 선이 그어진다.

\item[fitline] \keyvalueTF
이미지가 \macro{fitwidth}에 설정된 폭에 맞게 늘어나거나 줄어든다.

\item[fitwidth] \keyvalue{\linewidth}
이미지 폭

\item[flush] \keyvalueTF
변이단 조판에서 이미지가 왼쪽 여백으로 넘어간다.

\item[scale] \keyvalue{1.0}
주어진 배율만큼 이미지가 확대되거나 축소된다.

\item[landscape] \keyvalueTF
이미지가 90도 회전한다. 캡션도 함께 회전한다.

\item[middle] \keyvalueTF
이미지가 주어진 공간에서 세로로 중간에 놓인다.

\item[caption] \keyvalue{텍스트}
캡션 

\item[caption-short] \keyvalue{텍스트}
그림 차례에 들어가는 짧은 캡션

\item[caption-verb] \keyvalueTF
기호 문자들을 백슬래시 없이 그대로\annotate{verbatim} 캡션에 사용할 수 있다.

\item[caption-ignore] \keyvalueTF
캡션이 식자되지 않는다.

\item[caption-fit] \keyvalueTF
캡션의 폭이 이미지 폭에 맞춰진다.

\item[label] \keyvalue{텍스트}
상호참조를 위한 라벨. 달리 지정하지 않으면 이미지 파일 이름이 라벨이 된다.

\item[legend] \keyvalueTF
캡션에 번호가 붙지 않는다.

\item[showfilename] \keyvalueTF
이미지 파일의 이름이 표시된다.

\item[star] \keyvalue{옵션}
이것은 기본 설정과 다른 작동을 허용하기 위한 목적으로 고안되었다.
주어진 옵션이 별표(*)가 붙은 \macro{\image*}에 적용된다. 

\begin{code}
\ImageSetup{
    frame=false, scale=1, landscape=false, ...
    star={frame=true, scale=0.5, ...},
    vbar={landscape=true, ...}
}
\image{foo} \image*{foo} \image|{foo} \image*|{foo}
\end{code}

\item[vbar] \keyvalue{옵션}
\macro{star} 옵션과 마찬가지로, 주어진 옵션이 수직선(|)이 붙은 \macro{\image|}에 적용된다. 
\end{macros}

\macro{\listimg} 명령은 목록 환경에서 이미지를 삽입하기 위해 고안되었다. \macro{\image} 명령을 위한 전역 설정들 가운데 \macro{scale}과 \macro{showfilename} 옵션만이 이 명령에 적용된다.

\begin{coderesult}
\begin{enumerate}
\item 이 명령의 주요 목적은 소프트웨어 설명서에서 어떤 조작에 따르는 결과를 보여주는 것이다.
\listimg{uncertains}
\item 윗 글줄과 이미지 사이의 간격을 옵션으로 지정할 수 있다.
\listimg[-1ex]{uncertains}
\end{enumerate}
\end{coderesult}

\section{설명 딸린 이미지}

\macro{\illustimage} 명령 또는 \macro{IllustImage} 환경이 이미지와 설명문을 나란히 배치한다.

\begin{coderesult}
\begin{IllustImage}{uncertainm}
이 환경은 주어진 이미지에 대한 짧은 설명문을 곁에 두기 위한 목적으로 고안되었다.
\macro{wrapfigure} 환경과 달리, 글이 이미지 높이를 넘쳐도 이미지 아래로 흐르지 않는다.
\end{IllustImage}
\end{coderesult}

\begin{code}
\IllustImageSetup{옵션}
\illustimage*^[옵션]{이미지 파일}(캡션){텍스트}
\begin{IllustImage}*^[옵션]{이미지 파일}(캡션)
설명문
\end{IllustImage}
\end{code}

\begin{macros}<\IllustImageSetup>
\item[beforeskip] \keyvalue{0.25\onelineskip}
그림 위 간격

\item[extraskip] \keyvalue{0.5\onelineskip}
그림 위 간격. 
이 값이 삿갓표(\verb{^})가 붙은  \macro{\illustimage^} 또는 \macro{\begin{IllustImage}^} 에 적용된다. 

\item[afterskip] \keyvalue{0.25\onelineskip}
그림 아래 간격

\item[frame] \keyvalueTF
이미지 테두리에 선이 그어진다.

\item[gap] \keyvalue{1em}
이미지와 설명문 사이 간격

\item[scale] \keyvalue{1.0}
주어진 배율만큼 이미지가 확대되거나 축소된다.

\item[position] \keyvalue{left/right}
\macro{right}를 지정하면 설명이 왼쪽에, 이미지가 오른쪽에 배치된다.

\item[voffset] \keyvalue{0pt}
이미지 상단과 설명문 상단을 정렬하기 위해 설명문 위에 추가되는 간격

\item[valign] \keyvalue{t/c/b}
이미지와 설명문의 정렬 기준 (상단, 가운데, 하단)

\item[caption] \keyvalue{텍스트}
캡션 

\item[caption-verb] \keyvalueTF
기호 문자들을 백슬래시 없이 그대로\annotate{verbatim} 캡션에 사용할 수 있다.

\item[caption-ignore] \keyvalueTF
캡션이 식자되지 않는다.

\item[caption-align] \keyvalue{\raggedright, \centering, \raggedleft}
캡션 정렬

\item[label] \keyvalue{텍스트}
상호참조를 위한 라벨. 달리 지정하지 않으면 이미지 파일 이름이 라벨이 된다.

\item[legend] \keyvalueTF
캡션에 번호가 붙지 않는다.

\item[minwidth] \keyvalue{0pt}
이미지 영역의 최소 폭. 0 포인트보다 큰 값이 지정되었을 때 유효하다.

\item[imagealign] \keyvalue{\raggedright, \centering, \raggedleft}
이미지 영역 안에서 이미지 정렬

\item[broad] \keyvalueTF
변이단 조판에서 이미지 영역이 왼쪽 여백으로 넘어간다.

\item[showfilename]  \keyvalueTF
이미지 파일의 이름이 표시된다.

\item[textstyle] \keyvalue{\raggedright\small}
설명문 스타일 

\item[star] \keyvalue{옵션}
이것은 기본 설정과 다른 작동을 허용하기 위한 목적으로 고안되었다.
주어진 옵션이 별표가 붙은 \macro{\illustimage*} 또는 \macro{\begin{IllustImage}*} 에 적용된다. 

\item[enum] \keyvalue{옵션}
주어진 옵션이 \macro{\begin{IllustEnum}} 에 적용된다. 
\end{macros}

\macro{IllustEnum}은 설치 절차와 같은 순서를 나타내기 위한 목적으로 \macro{IllustImage}로부터 만들어진 환경이다.

\begin{coderesult}
\IllustImageSetup{enum=frame}
\begin{IllustEnum}*{uncertainm}
    \item 첫 단계를 나타내기 위해 별표를 붙인다.
\end{IllustEnum}
\begin{IllustEnum}^{uncertainm}
    \item 삿갓표를 붙이면 \macro{star}에 부여된 옵션들이 적용된다.
\end{IllustEnum}
\end{coderesult}

필요하다면 \macro{IllustImage}를 이용하여 다른 옵션들로 작동하는 환경들을 만들 수 있다.

\begin{code}
\tl_new:N \l_illust_item_options
\NewDocumentCommand \IllustItemSetup { m }
{
    \tl_set:Nn \l_illust_item_options { #1 }
}
\NewDocumentEnvironment { IllustItem }{ O{} m +b }
{
    \group_begin:
        \exp_args:No \IllustImageSetup{ \l_illust_item_options }
        \illustimage[#1]{#2}
        {
            \begin{itemize}
                #3
            \end{itemize}   
        }
    \group_end:
}{}
\end{code}

\section{글줄 안에 이미지}

스마트폰 설명서를 만드는 데에서 성가신 일들 가운에 하나가 버튼이나 아이콘 같은 작은 이미지들을 그려 넣은 것이다. \macro{\img} 명령이 이를 위해 고안되었다.

\begin{coderesult}
\LineImageSetup{scale=0.75}
이 버튼이 \img[-0.1ex]{button_menu} 흔히 메뉴를 가리킨다.
\end{coderesult}

\begin{code}
\LineImageSetup{옵션}
\image[높이]{이미지 파일}
\end{code}

\begin{macros}<\LineImageSetup>
\item[scale] \keyvalue{1.0}
주어진 배율만큼 이미지가 확대되거나 축소된다.

\item[raise] \keyvalue{-0.25ex}
이미지가 기준선\annotate{baseline}으로부터 주어진 크기만큼 올라가거나 내려간다.

\item[showfilename] \keyvalueTF
이미지 파일의 이름이 표시된다.
\end{macros}

\section{캡션}

\macro{\CaptionSetup} 명령을 이용하여 \term{캡션} 스타일을 바꿀 수 있다.

\begin{code}
\CaptionSetup{옵션}
\end{code}

\begin{macros}
\item[beforeskip] \keyvalue{1ex}
캡션 위 간격

\item[afterskip] \keyvalue{1ex}
캡션 아래 간격
\item[align] \keyvalue{\raggedright, \centering}
캡션이 두 줄 이상일 때 정렬 방식

\item[align-short] \keyvalue{\centering}
캡션이 한 줄 이하일 때 정렬 방식

\item[font] \keyvalue{\beseries\small}
폰트

\item[delimiter] \keyvalue{:}
번호와 캡션 사이 구획 문자
\end{macros}

\section{앨범 만들기}

텍으로 만들어진 문서를 인디자인 같은 다른 프로그램의 형식으로 옮겨야 할 때, 이미지 아래에 표시된 파일 이름이 이미지들을 식별하는 데에 도움이 된다.
\macro{\ShowImageFilename} 명령은 \macro{\ImageSetup}과 \macro{\IllustImageSetup}에, \macro{\ShowImageFilename*}은 \macro{\LineImageSetup}까지 포함하여, \macro{showfilename} 옵션을 지시한다.

\begin{code}
\MakeAlbum*[이미지 배율]{foo.txt}

foo.txt:
image_01.jpg
image_02.png
image_03.pdf
\end{code}

\macro{\MakeAlubm}은, 위 예와 같이, 텍스트 파일로부터 이미지 목록을 읽고 차례로 이미지들을 삽입한다.
\macro{\MakeAlbum*}은 이미지 파일의 이름을 식자하지 않는다.
다음 예는 현재 디렉토리에서 이미지 파일들을 모아 앨범을 만드는 간단한 \term{파워셸}\annotate{PowerShell} 스크립트이다.


\CodeSetup{language=powershell}
\begin{code}
$dir = "images.txt"
$tex = "album.tex"

$ImageTypes = @("pdf", "jpg", "png") 
foreach ($element in $ImageTypes) {
    Get-ChildItem "*.$element" -name | Add-Content $dir -Encoding UTF8
}

$src = "\documentclass{hzguide}
\LayoutSetup{paper=A4}
\begin{document}
\MakeAlbum{$dir} 
\end{document}"
Set-Content $tex -Encoding UTF8 $src

xelatex -interaction=batchmode $tex
\end{code}
\CodeSetup{language=latex}

\section{표}

\macro{Table} 환경이 표의 스타일을 제어한다.

\begin{code}
\TableSetup{옵션}
\begin{Table}*|[옵션](표 제목)<각주>
\begin{tabularx}{\linewidth}{lXX}
...
\end{tabularx}
\end{Table}
\end{code}

일반 출판물과 달리 설명서에서 표는 보충적인 자료가 아니라 핵심적인 정보의 일부이기 때문에 절 제목 아래에 오는 표 하나가, 다른 텍스트 없이도, 그 절의 온전한 내용이 될 수 있다. 따라서 불가피한 경우가 아니라면 표에 제목을 붙이지 않는 것이 바람직하다.

각주를 다는 방법에 대해 \pageref{sec:SpecTable} 페이지 \titleref{sec:SpecTable}\을 보라.

\begin{macros}<\TableSetup>
\item[beforeskip] \keyvalue{\hznulldim, 0.5\onelineskip}
표 위 간격. 값이 \macro{\hznulldim}이면 \macro{\ObjectSetup}으로 설정된 값이 사용된다.

\item[afterskip] \keyvalue{\hznulldim, 0.5\onelineskip}
표 아래 간격. \macro{\hznulldim}이면 \macro{\ObjectSetup}으로 설정된 값이 사용된다.

\item[float] \keyvalueTF
이 옵션은 표 내용을 \macro{table} 환경에 넣는다. 다시 말해 표의 위치가 텍에 의해 결정된다.

\item[float-placement] \keyvalue{tbh}
table 환경을 이용하는 경우에 위치 옵션

\item[align] \keyvalue{\raggedleft, \centering, \raggedright}
표 정렬

\item[landscape] \keyvalueTF
표가 90도 회전한다. 캡션도 함께 회전한다.

\item[caption] \keyvalue{텍스트}
캡션 

\item[label] \keyvalue{텍스트}
상호참조를 위한 라벨

\item[legend] \keyvalueTF
캡션에 번호가 붙지 않는다.

\item[font] \keyvalue{\rmfamily}
폰트 

\item[fontsize] \keyvalue{\small}
폰트 크기

\item[footnotestyle] \keyvalue{arabic, fnsymbol}
각주 번호 스타일

\item[stretch] \keyvalue{1.0}
줄 간격. 이것은 \macro{\arraystretch}를 설정하는 것과 같다.

\item[rowcolor] \keyvalueTF
줄 바탕이 번갈아 흰색과 회색이 된다.

\item[star] \keyvalue{옵션}
이것은 기본 설정과 다른 작동을 허용하기 위한 목적으로 고안되었다.
주어진 옵션이 별표가 붙은 \macro{\begin{Table}*}에 적용된다. 

\begin{code}
\TableSetup{
    rowcolor=true, fontsize=\small, ...,
    star={rowcolor=false, ...},
    vbar={font=\footnotesize, ...}
}
\begin{Table}
\begin{Table}*
\begin{Table}|
\begin{Table}*|
\end{code}

\item[vbar] \keyvalue{옵션}
\macro{star} 옵션과 마찬가지로, 주어진 옵션이 수직선이 붙은 \macro{\begin{Table}|}에 적용된다. 
\end{macros}

\subsection{제품 사양을 보여주는 표}
\label{sec:SpecTable}

\macro{SpecTable} 환경은 설명서에서 빈번하게 사용되는 표 형식 중 하나인 제품 사양을 보여주기 위한 장치이다. 
\macro{\TableSetup}에 의한 전역 설정이 이 환경에 유효하다.

\begin{code}
\begin{SpecTable}*[옵션](캡션)<각주>
...
\end{SpecTable}
\end{code}

\begin{coderesult}
\begin{SpecTable}[rowcolor]<
\footnotetext[1]{If the temperature inside the battery pack rises ...}
\footnotetext[2]{As the degradation of batteries is accelerated ...}
>
Nominal voltage\footnotemark[1] & 51.8 V \\
Operating voltage & 45.2 V to 58.1 V \\
Nominal capacity\footnotemark[2] & 126 A·h \\
Nominal energy & 6.4 kW·h \\
\end{SpecTable}
\end{coderesult}

별표 옵션(\macro{\begin{SpecTable}*})을 사용하면 첫째 칼럼에 \texttt{l} 형식이 적용된다.

\subsection{제품 구성품을 보여주는 표}

\macro{ImageTable} 환경은 제품과 함께 제공되는 액세서리를 보여주기 위한 장치이다.
\macro{\TableSetup}에 의한 전역 설정이 이 환경에도 적용된다.

\begin{code}
\begin{ImageTable}*<각주>
...
\end{ImageTable}
\end{code}

\begin{coderesult}
\begin{ImageTable}
\img{uncertains} & \img{uncertains} & \img{uncertains} \\
액세서리 & 액세서리 & 액세서리 \\
\end{ImageTable}
\end{coderesult}

별표 옵션(\macro{\begin{ImageTable}*})을 사용하면 칼럼 수가 둘로 줄어든다.

\chapter{목록과 용어}

\section{나열과 절차}

\macro{itemize}와 \macro{enumerate} 환경을 사용할 때 메므와 클래스가 제공하는 \macro{\tightlist} 명령이나 \macro{enumitem} 패키지가 제공하는 \macro{noitemsep} 옵션을 사용하여 항목 간격을 줄일 수 있다.
편의를 높이고자 별표를 사용하여 이를 지시할 수 있도록 \macro{itemize} 및 \macro{enumerate} 환경이 재정의되었다.
 
\begin{code}
\begin{enumerate}\tightlist
\begin{itemize}[noitemsep]
\begin{enumerate}*
\begin{itemize}*
\end{code}

필요한 경우가 드물겠지만, \macro{\SetListStyle} 명령을 이용하여 이 두 환경에 서로 다른 스타일을 적용할 수 있다.

\begin{code}
\SetListStyle{
    enumerate={\raggedright\small}
    itemize={\sffamily}
}
\end{code}

\section{정의 목록}

애플리케이션 프로그램의 설정 옵션들과 같은 것들을 설명하는 데에 \macro{description} 환경이 사용된다. 
\macro{enumitem} 패키지가 제공하는 \macro{\newlist} 명령을 이용하여 다른 형태의 정의 목록을 만들 수 있지만 흡족하지 않다.

\macro{\NewTerms} 명령은 다양한 형태의 정의 목록을 만드는 장치이다.

\begin{code}
\TermsSetup[환경 이름]{옵션}
\NewTerms{환경 이름}{옵션}
\RenewTerms{환경 이름}{옵션}
\end{code}

\macro{\TermsSetup}에 환경 이름을 지정하지 않으면 \macro{\NewTerms}에 의해 만들어진 모든 환경에 주어진 설정이 적용된다.

\begin{macros}<\TermsSetup>
\item[labeltype] \keyvalue{default, singleline, multiline, macro, macrosingle}
\macro{singleline}을 지정하면 표제어가 한 줄을 차지한다. 
\macro{macro}는 백슬래시를 비롯한 기호 문자들을 허용한다.
이 두 가지를 결합한 것이 \macro{macrosingle}이다.
\macro{multiline}을 지정하면, 표제어가 주어진 폭보다 길 때 두 줄이 된다.

\item[offset] \keyvalue{0em}
중첩된 환경에서 하위 수준의 왼쪽 여백에 추가되는 길이

\item[marker] \keyvalue{{{},\textbullet}}
표제어 앞에 붙이는 기호 문자. 앞의 것이 상위 수준에, 뒤의 것이 하위 수준에 적용된다.

\item[markersep] \keyvalue{0.5em}
기호 문자 뒤의 간격

\item[enumerate] \keyvalueTF
기호 문자 대신 번호가 붙는다.

\item[enummarker] \keyvalue{{\int_use:N \l_terms_int)\hskip\l_terms_marker_sep_skip}}
번호 형식

\item[base] \keyvalue{MM}
주어진 문자의 길이 만큼 표제어 영역의 폭이 정해진다.

\item[font] \keyvalue{\bfseries}
표제어 폰트

\item[color] \keyvalue{black}
표제어 색

\item[delimiter] \keyvalue{:\hfill}
표제어 뒤에 구획 문자

\item[labelwidth] \keyvalue{2em}
표제어 영역의 폭

\item[labelsep] \keyvalue{0.5em}
표제어 뒤의 간격

\item[itemindent] \keyvalue{0em}
표제어 뒤에 들여쓰기 폭

\item[leftmargin] \keyvalue{0pt}
왼쪽 여백

\item[index] \keyvalueTF
표제어가 색인에 추가된다.

\item[indexcategory] \keyvalue{단어}
표제어가 색인에서 주어진 단어 아래에 들어간다.

\item[indexsame] \keyvalueTF
표제어가 색인에서 지정된 것과 동일한 폰트로 식자된다.

\item[topsep] \keyvalue{1.25\onelineskip, 3ex}
목록 위 간격

\item[itemsep] \keyvalue{0.25\onelineskip, 1ex}
항목들 사이 간격

\item[style] \keyvalue{\raggedright, ...}
설명문 스타일

\item[beforelist] \keyvalue{\RuleBox, \FrameBox, ...}
환경 앞에 추가되는 명령. 
이를테면 이 옵션과 아래 옵션에 \macro{FrameBox} 환경의 시작 명령과 종료 명령을 지정함으로써 목록을 상자 안에 넣을 수 있다.

\item[afterlist] \keyvalue{\endRuleBox, \endFrameBox, ...}
환경 뒤에 추가되는 명령

\item[star] \keyvalue{옵션}
이것은 기본 설정과 다른 작동을 허용하기 위한 목적으로 고안되었다. 
주어진 옵션이 별표가 붙은 \verb{\begin{LIST}*}에 적용된다.

\item[vbar] \keyvalue{옵션}
\macro{star} 옵션과 마찬가지로, 주어진 옵션이 수직선이 붙은 \verb{\begin{LIST}|}에 적용된다.

\end{macros}

\macro{\NewTerms} 명령으로 다음과 같은 환경들이 저마다의 목적으로 정의되어 있다.

\begin{code}
\begin{terms}*|[옵션](표제어 폭)<상위 색인> \item[] \end{terms}
\begin{UI}*|[옵션](표제어 폭)<상위 색인> \item[] \end{UI}
\begin{macros}*|[옵션](표제어 폭)<상위 색인> \item[] \end{macros}
\begin{signs}*|[옵션](표제어 폭)<상위 색인> \item[] \end{signs}
\end{code}

\noindent "\verb{(MMMM)}"을 지정하는 것은 \verb{base=MMMM} 또는 \verb{labelwidth=4em}을 설정하는 것과 마찬가지이다.
"\verb{<단어>}"를  지정하는 것은 \verb{indexcategory=단어}를 설정하는 것과 같다.

\begin{macros}<\NewTerms>
\item[terms] 전문 용어\annotate{jargon}
\item[UI] 사용자 인터페이스 문자열
\item[macros] 레이텍 명령들 또는 그것들의 옵션.\footnote{이 문서의 작성을 돕기 위해 만들어졌다.} 표제어가 두 단어 이상으로 이루어져 있다면 \macro{\item[{}]} 또는 \macro{\vitem{}}을 이용하라. 공백을 포함하는 표제어들은 색인에 포함되지 않는다.
\item[signs] 안전 표지 같은 이미지. 이 환경에서 \verb{\item[]} 대신 \macro{\imgitem[]}을 사용해야 한다.
\end{macros}

\section{용어}
\label{sec:term}

\macro{\NewTerm} 명령을 이용하여 버튼 라벨 같은 말들을 강조하거나 색인에 넣는 명령들을 만들 수 있다.

\begin{code}
\TermSetup[명령 이름]{옵션}
\NewTerm{명령 이름}{옵션}
\RenewTerm{명령 이름}{옵션}
\end{code}

\macro{\TermSetup}에 명령 이름을 지정하지 않으면 \macro{\NewTerm}에 의해 만들어진 모든 명령에 주어진 설정이 적용된다.

\begin{macros}<\TermSetup>
\item[font] \keyvalue{\sffamily\bfseries}
용어 폰트

\item[color] \keyvalue{black}
용어 색

\item[index] \keyvalueTF
용어가 색인에 추가된다.

\item[indexsame] \keyvalueTF
표제어가 색인에서 지정된 것과 동일한 폰트로 식자된다.

\item[star] \keyvalue{옵션}
이것은 기본 설정과 다른 작동을 허용하기 위한 목적으로 고안되었다. 
주어진 옵션이 별표가 붙은 \verb{\TERM*}에 적용된다.

\item[vbar] \keyvalue{옵션}
\macro{star} 옵션과 마찬가지로, 주어진 옵션이 수직선이 붙은 \verb{\TERM|}에 적용된다.
\end{macros}

\macro{\NewTerm} 명령으로 다음과 같은 명령들이 저마다의 목적으로 정의되어 있다.

\begin{code}
\term*|[색인 정렬]{단어}[상위 색인]
\ui*|[색인 정렬]{단어}[상위 색인]
\mi*|[색인 정렬]{단어}[상위 색인]
\end{code}

\noindent \verb{\term[csharp]{C\#}[language]}는 \verb{\index{csharp@C\#}} 그리고 \verb{\index{language!csharp@C\#}}을 만들어낸다.

\begin{macros}<\NewTerm>
\item[\term] 전문 용어\annotate{jargon}
\item[\ui] 사용자 인터페이스 문자열
\item[\mi] 기기의 지시등이나 버튼에 표시된 문자열
\end{macros}

여러 \macro{\ui} 명령으로 이루어진 \macro{\menu} 명령은 메뉴 경로를 나타내기 위한 장치이다.

\begin{coderesult}
\menu[File][Preferences]{User Snippets}{latex.json}
\end{coderesult}

메므와 클래스에 정의된 \macro{\cmd} 명령을 이용하여 텍 명령을 식자할 수 있다.
이 명령은 색인 정렬을 위한 구분자로 "\verb{@}" 대신 "\verb{?}"를 사용한다.

\section{레이텍 명령}

\macro{\macro} 명령이 레이텍 명령들이나 그것들의 옵션을 고정폭 폰트로 식자한다.

\begin{code}
\macro{\includegraphics}
\macro{scale}
\end{code}

\section{첨자}

\macro{\textsuperscript}와 \macro{\textsubscript}의 다른 이름으로서 \macro{\Sup}와 \macro{\Sub}가 정의되어 있다.
화학식과 단위들은 정자로 표기되어야 하는데 수식을 이용하여 그것들을 표기하는 것이 도리어 번잡하기 때문이다.

\begin{coderesult}
$\mathrm{NH}_3$    NH\Sub{3} 
$\mathrm{μg/m}^3$  μg/m\Sup{3}
\end{coderesult}

\macro{\annotate} 명령은 \macro{\textsuperscript} 및 \macro{\textsubscript}를 대체하기 위해 고안되었다.

\begin{code}
\AnnoateSetup{옵션}
\annotate*[옵션]{주석 텍스트}
\anota*[옵션]{주석 텍스트}
\end{code}

\noindent \macro{\anota}는 \macro{\annotate}의 다른 이름이다.

\begin{macros}
\item[breakable] \keyvalueTF
이 옵션이 false이면 \verb{\textsuperscript}와 \verb{\textsubscript}가 사용된다. 이것은 줄 나눔이 허용되지 않는다는 것을 의미한다.

\item[superscript] \keyvalueTF
true이면 위 첨자로, false이면 아래 첨자로 식자된다.

\item[opening] \keyvalue{(}
주석 앞에 오는 문장 부호나 기호 문자

\item[closing] \keyvalue{)}
주석 뒤에 오는 문장 부호나 기호 문자

\item[font] \keyvalue{\scriptsize}
주석 폰트

\item[color] \keyvalue{darkgray}
주석 색

\item[space] \keyvalue{\thinspace}
주석 앞과 뒤에 오는 공백 문자

\item[raise] \keyvalue{0.75ex}
주석이 기준선으로부터 주어진 크기만큼 올라간다.
\macro{superscript} 옵션이 \false이면 내려간다.

\item[index] \keyvalueTF
주석이 색인에 추가된다.

\item[star] \keyvalue{옵션}
이것은 기본 설정과 다른 작동을 허용하기 위한 목적으로 고안되었다.
주어진 옵션이 별표가 붙은 \macro{\annotate*}에 적용된다.
\end{macros}

\section{색인}

HzGuide 클래스는, 다양한 언어들의 조판을 목적으로 하기 때문에, 색인 정렬에 \term{texindy}를 사용할 것을 권장한다.

\begin{code}
\IndexingEnable*
\IndexingDisable
\BookmarkIndexHead
\end{code}

\macro{\IndexingEnalbe}은 앞에서 설명한 색인을 지원하는 명령과 환경들의 \macro{index} 옵션을 \true로 변경한다.

\begin{itemize}*
\item \macro{\NewTerm} 명령에 의해 만들어진 명령들
\item \macro{\NewTerms} 명령에 의해 만들어진 환경들
\item \macro{\annotate} 명령
\item \macro{\macro} 명령
\end{itemize}

\macro{\BookmarkIndexHead} 명령은 색인 페이지에서 한글 자모 또는 알파벳을 가리키는 북마크를 PDF에 추가한다.
별표가 붙은 \macro{\IndexingEnable*} 명령이 \macro{\BookmarkIndexHead} 명령을 호출한다.

대부분의 경우에 \term{makeindex}나 \term{komkindex}가 문제를 일으키지 않을 것이다. 
그러나 \macro{\macro} 명령이나 \macro{macros} 환경을 이용하여 "\verb{@}"\annotate{at sign}가 포함된 레이텍 명령어를 색인에 넣고자 할 때에는 다음과 같은 모듈과 함께 \term{texindy} 또는 \term{xindy}를 사용해야 한다.

\begin{code}
(merge-rule "\macroprint{\(.*\)}" "\1~e" :again :bregexp)
(merge-rule "\termindexprint{.*}{\(.*\)}" "\1~e" :again :bregexp)
(merge-rule "?" "a~b")
\end{code}

\chapter{글상자와 경고문}

\section{글상자}

\macro{tcolorbox} 패키지에 기반하여 여러 모양의 박스들이 정의되어 있다.

\begin{FrameBox}
\begin{verbatim}
\begin{FrameBox} \end{FrameBox}
\end{verbatim}
\end{FrameBox} 
\bigskip

\begin{LeftBarBox}
\begin{verbatim}
\begin{LeftBarBox} \end{LeftBarBox}
\end{verbatim}
\end{LeftBarBox} 
\bigskip

\begin{RightBarBox}
\begin{verbatim}
\begin{RightBarBox} \end{RightBarBox}
\end{verbatim}
\end{RightBarBox} 
\bigskip

\begin{ShadeBox}[제목]
\begin{verbatim}
\begin{ShadeBox}[제목] \end{ShadeBox}
\end{verbatim}
\end{ShadeBox} 
\bigskip

\begin{ShadowBox}[제목]
\begin{verbatim}
\begin{ShadowBox}[제목] \end{ShadowBox}
\end{verbatim}
\end{ShadowBox} 
\bigskip

\begin{RuleBox}
\begin{verbatim}
\begin{RuleBox} \end{RuleBox}
\end{verbatim}
\end{RuleBox} 
\bigskip

\begin{TabBox}[제목]
\begin{verbatim}
\begin{TabBox}[제목] \end{TabBox}
\end{verbatim}
\end{TabBox}

이 박스들은 직접 사용되기보다, 특히 위와 아래 간격들이 0 포인트로 설정되어 있기 때문에, \macro{\NewBoundedBox} 명령이 정의하는 환경에서 사용될 목적으로 만들어졌다.

\begin{code}
\BoundedBoxSetup[환경 이름]{옵션}
\NewBoundedBox{환경 이름}{옵션}
\RenewBoundedBox{환경 이름}{옵션}
\end{code}

\macro{\BoundedBoxSetup}에 환경 이름을 지정하지 않으면 \macro{\NewBoundedBox}에 의해 만들어진 모든 환경에 주어진 설정이 적용된다.

\begin{macros}<\BoundedBoxSetup>
\item[beforeskip] \keyvalue{\hznulldim, 0.25\onelineskip}
박스 위 간격. 값이 \macro{\hznulldim}이면 \macro{\ObjectSetup}으로 설정된 값이 사용된다.

\item[extraskip] \keyvalue{-0.5\onelineskip}
그림 위 간격.
이 값이 삿갓표(\verb{^})가 붙은 \verb{\begin{BOX}^}에 적용된다. 

\item[afterskip] \keyvalue{\hznulldim, 0.25\onelineskip}
박스 아래 간격. 값이 \macro{\hznulldim}이면 \macro{\ObjectSetup}으로 설정된 값이 사용된다.

\item[type] \keyvalue{FrameBox, LeftBarBox, RightBarBox, ShadeBox, ShadowBox, RuleBox, TabBox}
박스 모양. 디폴트는 \macro{ShadeBox}이다.

\item[symbol] \keyvalue{\img{foo}\space}
제목 앞에 오는 이미지 

\item[signal] \keyvalue{표어}
특정한 목적을 위해 문서 전반에 걸쳐 사용되는 "경고"나 "주의" 같은 경고문 표어

\item[beforesignal] \keyvalue{<, \space}
제목 앞에 오는 문자

\item[aftersignal] \keyvalue{>, \space}
제목 뒤에 오는 문자

\item[style] \keyvalue{\raggedright}
문단 스타일

\item[titlefont] \keyvalue{\bfseries}
제목 폰트

\item[breakable-default] \keyvalueTF
이 옵션이 \true이면 박스의 쪽 나눔이 허용된다.

\item[breakable-adhoc] \keyvalueTF
이 옵션은 \macro{breakable-default}에 설정된 것과 반대로 설정되어야 한다.
이 설정이 별표가 붙은 \verb{\begin{BOX}*}에 적용된다.
\end{macros}

\macro{\NewBoundedBox} 명령으로 다음과 같은 환경들이 저마다의 목적으로 정의되어 있다.

심각성에 따른 경고문\annotate{admonition} 유형들은 \term{IEC/ISO 82079}나 \term{ASD-STE100} 같은 표준에 따라 약간 다르다.

\begin{danger}
\term{위험}은 사망이나 심각한 부상을 초래하는 상황을 가리킨다.
\macro{danger} 환경이 이렇게 정의되어 있다.
\begin{verbatim}
\NewBoundedBox{danger}{
    signal=\g_admonition_danger_tl, symbol={\img{AlertSymbol}~}
}
\end{verbatim}
\end{danger}

\begin{warning}
\term{경고}는 사망이나 심각한 부상을 초래할 수 있는 상황을 가리킨다.
\macro{warning} 환경이 이렇게 정의되어 있다.
\begin{verbatim}
\NewBoundedBox{warning}{
    signal=\g_admonition_warning_tl, symbol={\img{AlertSymbol}~}
}
\end{verbatim}
\end{warning}

\begin{caution}
\term{주의}\annotate{caution}는 경미한 부상한 초래하는 상황을 가리킨다.
\macro{caution} 환경이 이렇게 정의되어 있다.
\begin{verbatim}
\NewBoundedBox{caution}{signal=\g_admonition_caution_tl}
\end{verbatim}
\end{caution}

\begin{notice}
\term{주의}\annotate{notice}는 물적 피해를 초래하는 상황을 가리킨다.
\macro{notice} 환경이 이렇게 정의되어 있다.
\begin{verbatim}
\NewBoundedBox{notice}{signal=\g_admonition_notice_tl}
\end{verbatim}
\end{notice}

\begin{note}
\term{참고}는 주로 절차에서 추가 정보를 제공한다.
\macro{note} 환경이 이렇게 정의되어 있다.
\begin{verbatim}
\NewBoundedBox{note}{signal=\g_admonition_note_tl}
\end{verbatim}
\end{note}

\begin{reference}
\macro{reference} 환경은 파일 경로 같은 것을 예시하기 위한 것으로 이렇게 정의되어 있다.
\begin{verbatim}
\NewBoundedBox{reference}{style=\small}
\end{verbatim}
\end{reference}

\macro{\RenewTColorBox} 명령을 이용하여 이 박스들을 재정의할 수 있다.

\begin{code}
\RenewTColorBox{ShadeBox}{ s O{} }
{
    ... 
    title=#2, 
    IfBooleanTF={#1}{...}{...}
}
\end{code}

또한 다음과 같이 새로운 박스를 만들고 그 박스를 이용하는 환경을 만들 수 있다.

\begin{code}
\NewTColorBox{MyBox}{ s O{} }
{
    ...
}
\NewBoundedBox{MyBoxEnv}{
    type=MyBox,
    ...
}
\end{code}

\section{경고문 표어}

\macro{\RenewBoundedBox} 명령을 이용하여 \term{경고문 표어}\annotate{signal word}들을 변경하는 것이 번거롭기 때문에 \macro{\AdmonitionSetup} 명령이 고안되었다.
이를테면, \macro{language} 클래스 옵션에 \verb{german}이 지정되면 다음과 같은 선언이 발효된다.

\begin{code}
\AdmonitionSetup{
    danger=GEFAHR,
    warning=WARNUNG,
    caution=VORSICHT,
    notice=HINWEIS, 
    note=ANMERKUNG
}
\end{code}

\ChapterContentsDisable

\chapter{스타일 제어 도구}

\section{개체}
\label{sec:object}

\macro{\ObjectSetup}에 의한 설정이, 일괄적으로 제어할 수 있는 \term{개체}들로서, 다음 명령들과 환경들에 영향을 미친다.

\begin{itemize}
\item \macro{\image} 명령
\item \macro{IllustImage} 및 \macro{IllustEnum} 환경
\item \macro{callout} 환경
\item \macro{Table} 환경
\item \macro{caution}을 비롯한 \macro{\NewBoundedBox}에 의해 생성된 환경들
\item \macro{\action} 명령
\end{itemize}

\begin{macros}
\item[beforeskip] \keyvalue{0.5\onelineskip}
개체 위 간격.

\item[afterskip] \keyvalue{0.5\onelineskip}
개체 아래 간격

\item[framecolor] \keyvalue{gray}
테두리 색. 이를테면 \macro{\image} 명령의 \macro{frame} 옵션에 이 설정이 적용된다.

\item[align] \keyvalue{\raggedright}
정렬 방식
\end{macros}


\chapter{문서 표지}

\term{ISO/IEC Guide 37:2012}에 따르면 설명서는 제조사 이름, 주소, 웹 사이트를 비롯한 제조사 정보, 설명서가 적용되는 제품 유형과 모델을 비롯한 제품 정보, 설명서의 성격과 발행 날짜를 비롯한 문서 정보를 제공해야 한다. 
이 요건들을 충족하는 문서 표지를 \macro{\FrontCover}와 \macro{\BackCover} 명령이 생성한다.

\begin{code}
\CoverSetup{옵션}
\FrontCover
\BackCover
\end{code}

\begin{macros}
\item[topskip] \keyvalue{0.1\textheight}
앞표지에서 상단으로부터 띄우는 거리
\item[FrontLogoImage] \keyvalue{foo.jpg}
앞표지에 넣을 로고 이미지

\item[BackLogoImage] \keyvalue{foo.jpg}
뒤표지에 넣을 로고 이미지

\item[ProductImage] \keyvalue{foo.jpg}
제품 이미지

\item[FrontImage] \keyvalue{foo.jpg}
앞표지를 위한 배경 이미지. 
이미지를 놓는 기준점은 페이지 중앙이며, 페이지와 같은 크기의 이미지를 사용하는 것이 바람직하다.

\item[BackImage] \keyvalue{foo.jpg}
뒤표지를 위한 배경 이미지 

\item[title] \keyvalue{문서 제목}
"SM-G998" 같은 제품 모델

\item[subtitle] \keyvalue{부제}
"갤럭시 S21" 같은 제품 그룹 또는 유형.

\item[TitleAlign] \keyvalue{\raggedright}
제목 정렬

\item[DocumentType] \keyvalue{문서 유형}
"User Guide"나 "설치 설명서" 같은 문서의 목적 또는 대상 독자

\item[PubYear] \keyvalue{문서의 발행 년도}
달리 지정하지 않으면 현재 년도

\item[revision] \keyvalue{개정 번호}
이것은 일반 출판물의 판\annotate{版}과 같다.
제품의 수명 주기가 수년 이상일 때 그래서 여러 차례에 걸쳐 제품이 개선될 때 개정 번호가 유용할 수 있다. 
소프트웨어 제품의 경우에 그 버전을 나타내는 데에 이 옵션을 사용할 수도 있을 것이다.

\item[note] \keyvalue{문서에 대한 안내문}
가장 흔히 사용되는 문구는 "언제든지 필요할 때 볼 수 있도록 이 문서를 가까운 곳에 보관하십시오" 또는 "나중에 참조할 수 있도록 이 매뉴얼을 잘 보관하십시오"이다.

\item[manufacturer] \keyvalue{제조사의 공식 명칭}
"㈜데키스트" 또는 "Apple Inc."

\item[address] \keyvalue{제조사 주소}
본사 주소

\item[AfterFront] \keyvalue{\clearpage, \cleartooddpage}
앞표지 뒤에 쪽 나눔

\item[BeforeBack] \keyvalue{\clearpage, \cleartoevenpage}
뒤표지 앞에 쪽 나눔

\item[FontI] \keyvalue{\sffamily\bfseris}
제목과 부제를 위한 폰트

\item[FontII] \keyvalue{\sffamily\bfseris}
제목과 부제를 제외한 나머지에 적용되는 폰트

\item[FontSizeI] \keyvalue{\Huge}
제목 크기

\item[FontSizeII] \keyvalue{\huge}
부제 크기

\item[FontSizeIII] \keyvalue{\LARGE}
문서 유형의 크기

\item[FontSizeIV] \keyvalue{\Large}
발행 년도와 개정 번호의 크기

\item[BlankFront] \keyvalueTF
이 옵션이 \true로 설정되면, 배경 이미지를 제외하고 앞표지에 아무 것도 식자되지 않는다.
모든 문서 정보를 포함하는 이미지를 표지로 사용할 때 이 옵션이 유용하다.

\begin{code}
\CoverSetup{
    FrontImage=foo.jpg,
    BlankFront=true    
}
\end{code}

\item[BlankBack]
이 옵션이 \true로 설정되면, 배경 이미지를 제외하고 뒤표지에 아무 것도 식자되지 않는다.

\item[hook] \keyvalue{\foo}
이 옵션을 이용하여 모서리에 장식을 추가하는 것과 같은 명령을 삽입할 수 있다. 

\begin{code}
\usepackage[object=vectorian]{pgfornament}
\colorlet{OrnamentColor}{Maroon!60} 
\NewDocumentCommand \CoverOrnament { O{61} }
{	
\begin{tikzpicture}[every node/.style={inner sep=0pt}, remember picture, overlay]
\node[shift={(.5cm, -.5cm)}, anchor=north west, color=OrnamentColor] at (current page.north west) {\pgfornament[width=2cm]{#1}};
\node[shift={(-.5cm, -.5cm)}, anchor=north east, color=OrnamentColor] at (current page.north east) {\pgfornament[width=2cm, symmetry=v]{#1}};
\node[shift={(.5cm, .5cm)}, anchor=south west, color=OrnamentColor] at (current page.south west) {\pgfornament[width=2cm, symmetry=h]{#1}};
\node[shift={(-.5cm, .5cm)}, anchor=south east, color=OrnamentColor] at (current page.south east) {\pgfornament[width=2cm, symmetry=c]{#1}};
\end{tikzpicture}	
}
\CoverSetup{hook=\CoverOrnament}
\end{code}
\end{macros}

\printindex

\BackCover

\end{document}
